% !TEX program = xelatex
% Chinese layout in a book: secnum=zh chapters, fontsize=14pt, palatino,
% cjk-serif / cjk-serif-bold.
\documentclass[lang=zh, secnum=zh, fontsize=14pt, text=palatino,
  cjk-serif=FandolSong-Regular.otf, cjk-serif-bold=FandolSong-Bold.otf]{akizumi-book}
\title{論文}
\akizumisetup{subtitle={副標題}}
\begin{document}
\maketitle
\tableofcontents
\chapter{緒論}\label{ch:a}
\section{背景}\label{sec:a}
見\cref{ch:a}、\cref{sec:a}與\cref{thm:b}。\textbf{粗體}。
\begin{theorem}\label{thm:b} 內文。\end{theorem}
\ifdim\parindent=0pt \else\ERROR\fi  % par=skip stays the default
\makeatletter\ifdim\f@size pt=14pt \else\ERROR\fi\makeatother
\appendix
\chapter{證明}
\section{細節}
\end{document}
